\scanpage{41}\mspage{40}
% @title
\uline{Contrées induites}

% @ind
$1^\circ$) Soit $F$ une figure. Contrée \inserted{induite} $C_F$, \deleted{définie par les} c'est la partie fermée de $M_{\ll}$ engendrée par $F$, i.e. l'ens des $X^{\deleted{\unsure{F|f}}}\in M$ tels que $X$ raffine $F$.
% @lieu
% slanted
\margin{\uline{NB} L'ens des lieux de la contrée $C_F$ n'est autre que l'ombre de $F$}

% @resid
\uline{Contrée résiduelle} C'est la \inserted{sous} contrée $C_{\complement F}$ \deleted{\unsure{form|for}} formée des $X$ \inserted{$\in M$} qui sont disjoints de $F$. Par C\deleted{\unsure{o|on}}onst. 5, elle est \deleted{stable} \deleted{pour} fermée dans $M^{\unsure{\deleted{\leq}|\deleted{z}|\sharp}}_{\ll}$.
% @lieures
% slanted
\margin{\deleted{\unsure{Le|La}} L'ens des lieux de $C_{\complement F}$ est contenu dans $\complement_{\mathcal L}\widetilde F$, mais ne lui est \deleted{\unsure{égal|néces|pas}} pas \unsure{néc\textsuperscript{t}|néc\textsuperscript{t}.} égal.}

% @fig
\deleted{Cor} Soi\deleted{ent}t $F$ une figure, $G$ une \deleted{\uline{Subdivision}} sous figure. La contrée résiduelle de $F$ dans $G$ est par définition
% @meet
\[
 C_{\deleted{G}F}\cap C_{\complement\deleted{F}G}
\]

% @sub
\uline{Subdivision} Soit $F$ \inserted{$\subset M$} une figure, \deleted{et} On appelle \uline{subdivision} de $F$ une figure\unsure{\textsuperscript{o}|.} $F'\subset M$ telle que
% @a

a) $F'$ soit un raffinement de $F$
% @b

b) $F'$ soit maximale \inserted{pour $\subset$}, parmi les \deleted{raffi} figures qui raffinent $F$.
% @nb

\uline{NB} Si $F'$ raffine $F$, alors son ombre $\widetilde F'$ est contenue dans celle $\widetilde F$ de $F$.
